\section{Síntesis y ampliación}

\subsection{Repaso de la discusión central}
\begin{table}[H]
\centering
\begin{tabular}{p{2.8cm}p{4.0cm}p{5.0cm}}
\toprule
\textbf{Tema} & \textbf{Consenso del panel} & \textbf{Pregunta representativa} \\
\midrule
Contexto de la industria & 2025 es el año en que los Agent pasan de concepto a ciclo cerrado & ¿Por qué el capital, el producto y la investigación despegan al mismo tiempo? \\
Entrenamiento de modelos & Hay que construir el modelo como nativamente agentic, no depender solo de prompting & ¿Qué capacidades deberían interiorizarse en los parámetros del modelo? \\
Ingeniería de sistemas & El LM system sigue siendo clave; entre más fuerte el modelo, más necesita el sistema externo sostenerlo & ¿Cómo se diseñan memory, workflow y permisos? \\
Benchmark / RL & Un benchmark y un environment sólidos determinan si un método puede compararse de manera estable & ¿Cómo se construyen el reward y el environment del RL? \\
Dirección futura & Pasar de plantillas fijas a exploración autónoma, y seguir escalando acting y environment & ¿Dónde está el avance más importante para 2026? \\
\bottomrule
\end{tabular}
\caption{Comprime toda la discusión en cinco ejes de juicio}
\end{table}

\subsection{Tres conclusiones que vale la pena llevarse}
\begin{enumerate}
  \item La competencia entre Agent ya no es una competencia de capacidad de un solo modelo, sino una competencia conjunta entre capacidad del modelo, interfaz del sistema y construcción del environment.
  \item El cambio de investigación más importante de 2025 es empezar a tomar en serio la capacidad agentic como objeto de entrenamiento, y no solo como producto del prompt.
  \item Lo que de verdad limita la implementación a escala de los Agent no suele ser si el demo impresiona, sino si el environment, el benchmark, la memory, los permisos y el despliegue son lo bastante sólidos.
\end{enumerate}

\begin{importantbox}{Lo más valioso de este panel}
No cayó en la discusión superficial de «cuál framework es mejor», sino que trató al Agent como un sistema completo: hay que ver juntos el modelo, las herramientas, el environment, el entrenamiento, el producto y la implementación en la industria. Solo así «Agent» deja de ser un simple término de moda.
\end{importantbox}

\subsection{Lecturas adicionales}
\begin{itemize}
  \item Materiales públicos sobre sistemas de Agent de análisis de datos como Deep Research / Deep Analysis
  \item Artículos y evaluaciones relacionados con Mobile-Agent, GUI Agent y Computer Use
  \item Artículos sobre TinyZero, APR, multi-agent collaboration y la dirección de Agentic training
  \item Trabajos recientes sobre el uso de world model / environment model para el entrenamiento y la evaluación de Agent
  \item Prácticas de productos de Agent de código: materiales de diseño de flujos de trabajo sistemáticos de Cursor, Claude Code, etc.
\end{itemize}

\end{document}
